\documentclass[letterpaper]{article}

% ICML 2026 style
\usepackage{icml2026}

% Standard packages
\usepackage{amsmath,amssymb,amsfonts}
\usepackage{booktabs}
\usepackage{graphicx}
\usepackage{hyperref}
\usepackage{multirow}
\usepackage{xcolor}

\icmltitlerunning{Dynamic Health Networks: State-Conditional Dependency Modeling in Wearable Time Series}

\begin{document}

\twocolumn[
\icmltitle{Dynamic Health Networks: A Three-Level Framework for\\State-Conditional Dependency Modeling in Structured Wearable Time Series}

\icmlsetsymbol{equal}{*}

\begin{icmlauthorlist}
\icmlauthor{Anonymous}{}
\end{icmlauthorlist}

\vskip 0.3in
]

\printAffiliationsAndNotice{}

\begin{abstract}
We present a replicable three-level framework for modeling structured health time series as state-conditional dependency networks. Given daily multivariate physiological data from a consumer wearable device, the framework progressively builds (1) a lagged cross-correlation network capturing pairwise coupling strength and temporal precedence, (2) a Granger causality network adding directed, statistically controlled edges, and (3) a set of state-conditional networks that reveal how the dependency topology reorganizes across physiological contexts. We demonstrate the framework on 7 months of continuous Oura Ring data across six health dimensions (sleep, activity, recovery, blood oxygen, menstrual cycle, and exercise load). The state-conditional analysis uncovers a hub rotation phenomenon: sleep anchors the baseline network, activity dominates during the luteal phase, and recovery becomes central after competitive tennis matches. We further identify a sign reversal in the sleep--activity relationship and a Granger causality direction reversal between cycle phases, providing observational causal evidence for a mode shift from volitional to capacity-limited physiology. The framework is released as open-source code (\url{https://github.com/builderfeng-playera/dynamic-health-networks}) applicable to any consumer wearable dataset, requiring no clinical infrastructure.
\end{abstract}

%----------------------------------------------------------------------
\section{Introduction}
%----------------------------------------------------------------------

Consumer wearable devices now generate continuous, multi-dimensional structured health data, including sleep architecture, heart rate variability (HRV), activity patterns, blood oxygen, skin temperature, and stress indicators, enabling longitudinal health monitoring at unprecedented granularity. However, the dominant analytical paradigm treats these metrics independently: tracking individual signals against personal baselines \citep{mishra2020}, validating sensor accuracy per metric \citep{cao2022}, or computing summary statistics across observation periods \citep{huhn2022}. When inter-metric relationships are examined, they are typically modeled as fixed pairwise correlations over the entire dataset. This static view obscures a fundamental property of physiological systems: \textbf{the dependency structure among health metrics is itself state-dependent.}

Network Physiology has established that organ-level interactions dynamically reorganize across physiological states. \citet{bartsch2015} demonstrated that brain--heart--respiration coupling networks undergo pronounced topological changes across sleep stages. \citet{ivanov2021} formalized this into a broader vision of the ``Human Physiolome,'' the complete map of dynamic organ-to-organ interaction networks, and called for extending this framework to consumer wearable data. Yet no study has provided a replicable, end-to-end analytical framework that takes consumer-grade structured health time series as input and outputs state-conditional dependency networks with formal significance testing.

We address this gap with a three-level framework for structured health time series modeling. The framework progressively constructs: (L1) lagged cross-correlation networks capturing pairwise coupling and temporal precedence, (L2) Granger causality networks adding directed edges with causal interpretation from observational data, and (L3) state-conditional networks that partition the time series by physiological context and quantify topological reorganization via network density, hub centrality, and Frobenius divergence. Each level builds on the previous, and the entire pipeline is released as open-source code.\footnote{\url{https://github.com/builderfeng-playera/dynamic-health-networks}}

We demonstrate the framework on 7 months of Oura Ring data across six health dimensions (sleep, activity, recovery, blood oxygen, menstrual cycle, and exercise load), showing that the network topology reorganizes across biological and behavioral states, with the \textbf{hub node rotating} depending on context. Notably, the Granger causality analysis reveals a direction reversal between cycle phases, providing observational causal evidence for a physiological mode shift. These findings also highlight that current wearable platforms do not condition metric interpretation on menstrual cycle state \citep{alzueta2022,goodale2019}, and sport science still lacks standardized cycle-aware monitoring protocols \citep{carmichael2025}.

\textbf{Contributions.} (1) A replicable, open-source three-level framework for modeling structured health time series as state-conditional dependency networks. (2) A demonstration that Granger causality applied to consumer wearable data can recover directed, state-dependent physiological relationships. (3) Empirical evidence that state-conditional analysis reveals network dynamics (hub rotation, sign reversals, density shifts) invisible to static correlation methods.

%----------------------------------------------------------------------
\section{Data \& Setup}
%----------------------------------------------------------------------

\subsection{Data Source}

We analyze 7 months of continuous data (September 8, 2025 -- April 13, 2026) collected from an Oura Ring Generation 3 worn by a single healthy female participant (age 33). The ring captures sleep architecture, nocturnal heart rate and HRV, daily activity, skin temperature, blood oxygen saturation, and autonomic stress/recovery states via photoplethysmography (PPG), accelerometer, and NTC thermistor sensors. Of 217 calendar days, 168 are classified as usable after excluding a 29-day data gap (travel without charger), a 14-day cold-start calibration period, and 7 low-wear days. The dataset is not publicly released; our contribution is the analytical method, not the data.

\subsection{Network Nodes}

We model the body as a 6-node network, where each node represents a physiological dimension aggregated from multiple raw metrics. Three nodes use Oura's built-in composite scores, proprietary weighted aggregates computed by the device firmware and presented to all Oura users globally. Three nodes without built-in scores are constructed from raw metrics. Table~\ref{tab:nodes} summarizes the node definitions.

\begin{table}[t]
\caption{Network node definitions.}
\label{tab:nodes}
\vskip 0.1in
\centering
\small
\begin{tabular}{@{}clp{4.2cm}@{}}
\toprule
Node & Dimension & Construction \\
\midrule
\textbf{A} & Sleep quality & Oura sleep score (0--100) \\
\textbf{B} & Daily activity & Oura activity score (0--100) \\
\textbf{C} & Recovery/stress & Oura readiness score (0--100) \\
\textbf{D} & Blood oxygen & Mean of z-scored SpO\textsubscript{2} and inverted breathing disturbance index \\
\textbf{E} & Menstrual cycle & Skin temperature deviation from baseline (\textdegree C) \\
\textbf{F} & Tennis load & Composite of match indicator, workout peak HR (z), active calories (z) \\
\bottomrule
\end{tabular}
\end{table}

\subsection{Body States}

We partition the 168 usable days into three physiological states based on biological markers and known events (Table~\ref{tab:states}). States are mutually exclusive; days falling into multiple categories are assigned by priority: post-tennis match $>$ luteal $>$ baseline. Days surrounding the 29-day data gap (travel absence, Feb 26 -- Mar 26) are excluded due to insufficient sample size for reliable network estimation.

\begin{table}[t]
\caption{Body state definitions.}
\label{tab:states}
\vskip 0.1in
\centering
\small
\begin{tabular}{@{}lcp{3.5cm}@{}}
\toprule
State & $N$ days & Identification \\
\midrule
Baseline (Follicular) & $\sim$70 & Temp.\ deviation $\leq$ 0\textdegree C, no match flag \\
Luteal phase & $\sim$55 & Temp.\ deviation $>$ 0\textdegree C sustained $\geq$ 3 days \\
Post-tennis match & 8 & Match day $+$ 1 day after (4 matches $\times$ 2) \\
\bottomrule
\end{tabular}
\end{table}

%----------------------------------------------------------------------
\section{Methods}
%----------------------------------------------------------------------

We analyze inter-node dependencies at three levels of increasing complexity. Each level builds on the previous, progressively adding directionality and state-conditioning.

\subsection{Level 1: Lagged Cross-Correlation Network}

For each pair of nodes $(X_i, X_j)$, we compute the Pearson cross-correlation at lags $\tau = 0, 1, \ldots, 7$ days:
\begin{equation}
r_{ij}(\tau) = \text{Corr}\bigl(X_i(t),\; X_j(t + \tau)\bigr)
\end{equation}
The optimal lag $\tau^*_{ij} = \arg\max_\tau |r_{ij}(\tau)|$ identifies the strongest coupling and its temporal offset. An edge is drawn if $|r_{ij}(\tau^*)| > r_{\text{thresh}}$, where $r_{\text{thresh}}$ is the 95th percentile of $|r|$ under 1{,}000 permutations of one series.

\subsection{Level 2: Granger Causality Network}

To test whether the past of $X_i$ improves prediction of $X_j$ beyond $X_j$'s own past, we fit two autoregressive models of order $p$:
\begin{align}
\text{Restricted:}\quad X_j(t) &= \textstyle\sum_{k=1}^{p} \alpha_k\, X_j(t\!-\!k) + \epsilon_t \\
\text{Unrestricted:}\quad X_j(t) &= \textstyle\sum_{k=1}^{p} \alpha_k\, X_j(t\!-\!k) + \textstyle\sum_{k=1}^{p} \beta_k\, X_i(t\!-\!k) + \eta_t
\end{align}
A standard $F$-test compares residual sums of squares ($p < 0.05$, Bonferroni-corrected for 30 pairwise tests). Model order $p$ is selected by BIC over $p \in \{1, \ldots, 7\}$. Stationarity is ensured by first-differencing any series that fails the augmented Dickey--Fuller test. For states with fewer than 30 days, we report only L1 results.

\subsection{Level 3: State-Conditional Dynamic Network}

We compute separate networks for each body state $s \in \{\text{baseline, luteal, post-match}\}$. For each state, we extract the corresponding subsequences and compute the pairwise correlation matrix $\mathbf{C}^{(s)} \in \mathbb{R}^{6 \times 6}$. Edge significance is assessed via block bootstrap (contiguous blocks of 3 days, 2{,}000 iterations); an edge is retained if the 95\% confidence interval for $|r_{ij}|$ excludes zero.

\textbf{Network comparison metrics.} We quantify structural differences using three measures:

\begin{enumerate}
\item \textbf{Network density} $\rho^{(s)}$: fraction of 15 possible edges that are significant.
\item \textbf{Hub centrality} $h^{(s)}_i$: degree centrality of node $i$ (significant edges / 5). The dominant hub is $\arg\max_i h^{(s)}_i$.
\item \textbf{Network divergence} $\Delta^{(s)} = \|\mathbf{C}^{(s)} - \mathbf{C}^{(\text{baseline})}\|_F$: Frobenius distance from the baseline topology.
\end{enumerate}

%----------------------------------------------------------------------
\section{Results}
%----------------------------------------------------------------------

\subsection{Global Network Structure (L1)}

Across all 168 usable days, the lag cross-correlation network reveals 4 significant edges out of 15 possible (density $= 0.27$, permutation threshold $|r| > 0.217$). The strongest connections are: Activity--Cycle ($r = +0.420$, lag 1d), Activity--Tennis ($r = +0.337$, lag 0d), Recovery--Cycle ($r = -0.312$, lag 0d), and Recovery--Tennis ($r = +0.219$, lag 6d). Notably, the sleep node (A) has no significant global edges; its connections are state-dependent and wash out in the aggregate, motivating the state-conditional analysis.

\subsection{Directed Influences (L2)}

Granger causality analysis on the two states with sufficient sample size (baseline: 65 days, luteal: 81 days) reveals sparse but interpretable directed edges. In baseline, Recovery Granger-causes Blood Oxygen ($\text{C} \to \text{D}$, $F = 9.21$, $p = 0.004$), and Tennis Granger-causes Activity ($\text{F} \to \text{B}$, $F = 4.26$, $p = 0.019$). In the luteal state, the directed structure shifts: Activity Granger-causes Tennis ($\text{B} \to \text{F}$, $F = 4.41$, $p = 0.039$), \textbf{reversing} the baseline direction. This reversal suggests that during the follicular phase, exercise choice drives activity, while during the luteal phase, the body's general capacity constrains exercise output, consistent with progesterone-mediated fatigue \citep{mcnulty2020}.

\subsection{State-Conditional Networks (L3)}

Constructing separate networks for each body state reveals pronounced topological reorganization (Table~\ref{tab:results}).

\begin{table}[t]
\caption{State-conditional network summary.}
\label{tab:results}
\vskip 0.1in
\centering
\small
\begin{tabular}{@{}lcccr@{}}
\toprule
State & $N$ & $\rho$ & Dom.\ Hub & $\Delta^{(s)}$ \\
\midrule
Baseline (Follicular) & 65 & 0.20 & A (Sleep) & -- \\
Luteal phase & 81 & 0.20 & B (Activity) & 1.07 \\
Post-tennis match & 8 & \textbf{0.40} & \textbf{C (Recovery)} & \textbf{2.42} \\
\bottomrule
\end{tabular}
\end{table}

\textbf{Hub rotation.} The dominant hub shifts across states: Sleep (A) anchors the baseline, Activity (B) takes over during the luteal phase, and Recovery (C) dominates post-tennis. No single node is universally central; the body's coordination center rotates depending on what physiological challenge is active.

\textbf{Post-tennis densification.} The post-tennis network is the most interconnected (density 0.40, double baseline) with the largest divergence ($\Delta = 2.42$). Three edges involving Recovery and Tennis emerge with strong magnitudes: C--F ($r = +0.848$), A--F ($r = +0.778$), and C--E ($r = -0.473$). The appearance of a Cycle--Recovery edge absent in baseline suggests that menstrual cycle phase modulates recovery capacity even during acute exercise recovery, a cross-system interaction invisible in static analysis.

\subsection{Edge-Level State Transitions}

Table~\ref{tab:edges} shows how specific edges change across states, revealing the mechanism behind hub rotation.

\begin{table}[t]
\caption{Key edge transitions across body states. $\checkmark$ = significant (bootstrap $p < 0.05$).}
\label{tab:edges}
\vskip 0.1in
\centering
\small
\begin{tabular}{@{}lccc@{}}
\toprule
Edge & Baseline & Luteal & Post-Tennis \\
\midrule
A--B (Sleep--Act.) & $+0.25$ $\checkmark$ & $-0.17$ $\checkmark$ & n.s. \\
A--C (Sleep--Rec.) & $+0.27$ $\checkmark$ & n.s. & $+0.66$ $\checkmark$ \\
B--F (Act.--Tennis) & $+0.38$ $\checkmark$ & $+0.47$ $\checkmark$ & n.s. \\
C--E (Rec.--Cycle) & n.s. & $-0.31$ $\checkmark$ & $-0.47$ $\checkmark$ \\
C--F (Rec.--Tennis) & n.s. & n.s. & $+0.85$ $\checkmark$ \\
\bottomrule
\end{tabular}
\end{table}

The Sleep--Activity sign reversal (A--B: $+0.245 \to -0.171$) between follicular and luteal phases is particularly notable. During the follicular phase, better sleep aligns with higher activity. During the luteal phase, the relationship inverts: the body trades off between sleep quality and activity output, consistent with progesterone-mediated fatigue competing with daytime demands.

%----------------------------------------------------------------------
\section{Discussion}
%----------------------------------------------------------------------

\subsection{Physiological Interpretation}

\textbf{Luteal phase.} Progesterone increases basal metabolic rate, elevates core temperature by 0.2--0.5\textdegree C, and shifts autonomic balance toward sympathetic dominance \citep{baker2007}. This explains the new inhibitory Cycle--Recovery edge ($r = -0.31$). The simultaneous sign reversal in Sleep--Activity and the Granger causality direction reversal together suggest a mode shift from volitional to capacity-limited physiology, where the body constrains what activities are feasible.

\textbf{Post-tennis.} Acute high-intensity exercise triggers elevated cortisol, suppressed HRV, glycogen depletion, and inflammatory signaling \citep{halson2014}. The network densification and Recovery emerging as dominant hub reflect whole-body mobilization. The exceptionally strong C--F edge ($r = +0.85$) indicates tight coupling between exercise load and recovery demand, absent in baseline. The Cycle--Recovery edge also activating post-match suggests that cycle phase modulates recovery capacity even under acute exercise stress.

\subsection{Framework Applicability}

The three-level framework is portable across consumer wearable platforms and user populations. The input is a daily multivariate time series with a state partition; the output is a set of state-conditional networks with formal significance testing. The node construction step is the only component requiring device-specific adaptation; the analytical levels (L1--L3) operate on any set of continuous daily time series. We release the full pipeline as open-source code.\footnote{\url{https://github.com/builderfeng-playera/dynamic-health-networks}}

The framework extends naturally to richer state partitions. In this study we used three states defined by biological markers; in principle, any contextual annotation (shift work schedules, medication changes, seasonal variation) can serve as the conditioning variable. This makes the framework a general-purpose tool for structured health time series analysis.

\subsection{Limitations}

\textbf{Sample size.} This is an N-of-1 study with 168 usable days. The post-tennis state (8 days) has a limited sample, reducing statistical power and precluding Granger analysis. Effects in this state should be interpreted as hypothesis-generating.

\textbf{Node construction.} Three nodes rely on Oura's proprietary composite scores, whose internal weighting is not publicly documented, introducing a dependency on vendor-specific algorithms.

\textbf{State partitioning.} The luteal/follicular classification uses temperature deviation as a proxy. Without hormonal assays, exact ovulation timing is approximate.

\textbf{Causality.} Granger causality tests temporal precedence, not true causal mechanisms. Directed edges should be interpreted as predictive relationships.

\textbf{Equity considerations.} Current wearable platforms report readiness and recovery scores without conditioning on cycle phase, potentially introducing systematic bias for female users. Extending the framework to diverse populations would help quantify such biases.

\subsection{Future Work}

The framework can be extended with (L4) transfer entropy for non-linear directional information flow \citep{schreiber2000}, and (L5) regime-switching VAR models that infer body states from data rather than predefined partitions \citep{hamilton1989}. The immediate next step is multi-subject validation to test whether hub rotation generalizes across individuals.

%----------------------------------------------------------------------
% References
%----------------------------------------------------------------------

\bibliography{references}
\bibliographystyle{icml2026}

\end{document}
